利用 Green 函数和紧算子理论, 可以证明以下结果.

\begin{Theorem}
\label{thm:regular-sturm-liouville-complete-spectral}
设 $(\mathcal{D}(L),L)$ 由 \eqref{eq:sturm-liouville-weighted-differential-operator} 和 \eqref{eq:sturm-liouville-unmixed-operator-domain} 定义. 以下结论成立:\begin{enumerate}\item 谱完全由特征值组成.\item 特征值可数, 并可列成序列
\begin{equation}\label{eq:sturm-liouville-increasing-eigenvalue-sequence}
\lambda_1<\lambda_2<\cdots<\lambda_n<\cdots
\end{equation}
且
\begin{equation}\label{eq:sturm-liouville-eigenvalues-tend-to-infinity}
\lim_{n\to\infty}\lambda_n=\infty.
\end{equation}
\item 满足 $\phi_i\in\mathcal{N}(L_{\lambda_i})$ 的归一化特征函数集 $\{\phi_i\}$ (其中 $L_\lambda:=L-\lambda I$) 是 $L_w^2(a,b)$ 的一组标准正交基.\item 对于方程
\begin{equation}\label{eq:sturm-liouville-spectral-inhomogeneous-equation}
Lu-\lambda u=f,
\end{equation}
以下两种情形恰有一种成立:\begin{enumerate}[label=(\alph*)]\item 若 $\lambda$ 不是 $(\mathcal{D}(L),L)$ 的特征值, 则对每个 $f\in L_w^2(a,b)$, \eqref{eq:sturm-liouville-spectral-inhomogeneous-equation} 在 $\mathcal{D}(L)$ 中有唯一解. 该解为
\begin{equation}\label{eq:sturm-liouville-noneigenvalue-solution-expansion}
u(x)=\sum_{i=1}^\infty\dfrac{(\phi_i,f)_w}{\lambda_i-\lambda}\phi_i(x).
\end{equation}
\item 若 $\lambda=\lambda_j$ 是 $(\mathcal{D}(L),L)$ 的特征值, 则当
\begin{equation}\label{eq:sturm-liouville-eigenvalue-solvability-condition}
(\phi_j,f)=0.
\end{equation}
时, \eqref{eq:sturm-liouville-spectral-inhomogeneous-equation} 在 $\mathcal{D}(L)$ 中有解. 此时存在如下单参数解族:
\begin{equation}\label{eq:sturm-liouville-eigenvalue-solution-family}
u(x)=C\phi_j(x)+\sum_{\substack{i=1\\i\ne j}}^\infty\dfrac{(\phi_i,f)_w}{\lambda_i-\lambda}\phi_i(x).
\end{equation}
\end{enumerate}\end{enumerate}
\end{Theorem}

\begin{Proof}
首先为 ODE
\begin{equation}\label{eq:sturm-liouville-homogeneous-expanded-ode}
-pu''-p'u'+(q-\lambda w)u=0
\end{equation}
及边界条件 \eqref{eq:sturm-liouville-left-unmixed-boundary-condition} 和 \eqref{eq:sturm-liouville-right-unmixed-boundary-condition} 构造 Green 函数. 设 $v_l(x;\lambda)$ 满足 ODE \eqref{eq:sturm-liouville-homogeneous-expanded-ode} 和左端边界条件 \eqref{eq:sturm-liouville-left-unmixed-boundary-condition}, $v_r(x;\lambda)$ 满足 \eqref{eq:sturm-liouville-homogeneous-expanded-ode} 和右端边界条件. 假设 $\lambda$ 不是特征值. 此时, ODE 初值问题的唯一性定理说明 $v_l$ 和 $v_r$ 线性无关.

Green 函数将具有以下形式:
\begin{equation}\label{eq:sturm-liouville-green-function-piecewise-ansatz}
g(x,y;\lambda):=\begin{cases}
a_l(y;\lambda)v_l(x;\lambda),&a<x<y,\\
a_r(y;\lambda)v_r(x;\lambda),&y<x<b.
\end{cases}
\end{equation}
于是 $g(\cdot,y;\lambda)$ 满足 \eqref{eq:sturm-liouville-left-unmixed-boundary-condition} 和 \eqref{eq:sturm-liouville-right-unmixed-boundary-condition}. 为得到
\begin{equation}\label{eq:sturm-liouville-green-function-delta-equation}
-p\frac{\partial^2}{\partial x^2}g(x,y;\lambda)-p'\frac{\partial}{\partial x}g(x,y;\lambda)+(q-\lambda w)g(x,y;\lambda)=\delta(x-y)
\end{equation}
要求\[\begin{aligned}
a_l(y;\lambda)v_l(y;\lambda)-a_r(y;\lambda)v_r(y;\lambda)&=0,\\
a_l(y;\lambda)v_l'(y;\lambda)-a_r(y;\lambda)v_r'(y;\lambda)&=\dfrac{1}{p(y)}.
\end{aligned}\]解得
\begin{equation}\label{eq:sturm-liouville-green-left-coefficient}
a_l(y;\lambda)=(W(v_l(y;\lambda),v_r(y;\lambda))p(y))^{-1}v_r(y;\lambda),
\end{equation}

\begin{equation}\label{eq:sturm-liouville-green-right-coefficient}
a_r(y;\lambda)=(W(v_l(y;\lambda),v_r(y;\lambda))p(y))^{-1}v_l(y;\lambda),
\end{equation}
其中
\begin{equation}\label{eq:sturm-liouville-fundamental-wronskian}
W(v_l(y;\lambda),v_r(y;\lambda)):=v_l(y;\lambda)v_r'(y;\lambda)-v_l'(y;\lambda)v_r(y;\lambda)
\end{equation}
为 $v_l$ 和 $v_r$ 的 Wronskian 行列式. 一个称为 \emph{Abel 公式} 的经典 ODE 结果 (这里不予证明) 给出
\begin{equation}\label{eq:sturm-liouville-abel-constant-wronskian}
W(v_l(y;\lambda),v_r(y;\lambda))p(y)=C^{-1}
\end{equation}
其中 $C$ 为常数. 因而有
\begin{equation}\label{eq:sturm-liouville-green-function-explicit-kernel}
g(x,y;\lambda):=\begin{cases}
Cv_r(y;\lambda)v_l(x;\lambda),&a<x<y,\\
Cv_l(y;\lambda)v_r(x;\lambda),&y<x<b.
\end{cases}
\end{equation}
注意, $g$ 有界, 而且由于 $w(x)$ 有界, 核 $w(y)g(x,y;\lambda)$ 为 Hilbert-Schmidt 核. 因而算子
\begin{equation}\label{eq:sturm-liouville-weighted-green-resolvent-operator}
L_w^2(a,b)\ni u\mapsto K_\lambda u:=\int_a^bw(y)g(\cdot,y;\lambda)u(y)\,\mathrm{d}y\in L_w^2(a,b)
\end{equation}
有界. 直接计算表明 $K_\lambda=R_\lambda(L)$. 由于 $K_\lambda$ 的定义域稠密, 且当且仅当 $\lambda$ 不是 $(\mathcal{D}(L),L)$ 的特征值时才有定义, 第 1 条得证.

由前一定理的第 2 条, 可不失一般性地假设 $\lambda=0$ 不是特征值 (否则适当改变 $q$). 令 $g_0(x,y):=g(x,y;0)$. 注意, 积分方程
\begin{equation}\label{eq:sturm-liouville-kernel-eigenproblem}
u(x)=\lambda\int_a^bw(y)g_0(x,y)u(y)\,\mathrm{d}y
\end{equation}
存在非平凡解 $u$, 当且仅当 $\lambda$ 为 $(\mathcal{D}(L),L)$ 的特征值, 且 $u$ 为相应特征函数. 令\[\begin{aligned}
v(x)&:=\sqrt{w(x)}u(x),\\
k(x,y)&:=\sqrt{w(x)}g_0(x,y)\sqrt{w(y)},
\end{aligned}\]即可将 \eqref{eq:sturm-liouville-kernel-eigenproblem} 写为
\begin{equation}\label{eq:sturm-liouville-unweighted-compact-eigenproblem}
Gv(x):=\int_a^bk(x,y)v(y)\,\mathrm{d}y=\dfrac{1}{\lambda}v(x).
\end{equation}
由于 $k$ 对称且有界, $G$ 为从 $L^2(a,b)$ 到 $L^2(a,b)$ 的自伴随紧算子. 显然, $(\mathcal{D}(L),L)$ 的特征值为 $G$ 的特征值的倒数, 而 $L$ 的归一化特征函数可由 $G$ 的归一化特征函数通过公式 $\phi_i(x)=w(x)^{-1/2}v_i(x)$ 得到. 此外, 由于对任意 $f\in L_w^2(a,b)$ 均有 $\sqrt{w}f\in L^2(a,b)$, 所以可写为
\begin{equation}\label{eq:sturm-liouville-weighted-data-basis-expansion}
\sqrt{w(x)}f(x)=\sum_{i=1}^\infty(v_i,\sqrt{w}f)v_i(x)=\sqrt{w(x)}\sum_{i=1}^\infty(\phi_i,f)_w\phi_i(x).
\end{equation}
于是
\begin{equation}\label{eq:sturm-liouville-data-eigenfunction-expansion}
f(x)=\sum_{i=1}^\infty(\phi_i,f)_w\phi_i(x).
\end{equation}
第 2 和第 3 条得证.

非齐次方程解的存在性和重数由 Fredholm 二择一定理及等价的积分形式得到. 为导出展开公式, 先假设 $\lambda$ 不是特征值. 令
\begin{equation}\label{eq:sturm-liouville-solution-basis-expansion}
u(x)=\sum_{i=1}^\infty(\phi_i,u)_w\phi_i(x)
\end{equation}
且
\begin{equation}\label{eq:sturm-liouville-forcing-basis-expansion}
f(x)=\sum_{i=1}^\infty(\phi_i,f)_w\phi_i(x).
\end{equation}
则
\begin{equation}\label{eq:sturm-liouville-equating-spectral-coefficients}
(L-\lambda I)u=\sum_{i=1}^\infty(\phi_i,u)_w(\lambda_i-\lambda)\phi_i(x)=\sum_{i=1}^\infty(\phi_i,f)_w\phi_i(x).
\end{equation}
比较系数即得 \eqref{eq:sturm-liouville-noneigenvalue-solution-expansion}. 类似计算给出 \eqref{eq:sturm-liouville-eigenvalue-solution-family}. 证明完毕.
\end{Proof}
